%% ================================================================
%%  SECTION 10 — Gaps and Contribution
%% ================================================================
\section{Gaps and Contribution}

We can state the central finding of this survey plainly: nobody has built it yet. No existing Indian legal AI system connects precedent retrieval to case fact structuring, structuring to argument graph construction, argument graphs to outcome prediction, or prediction to counterfactual what-if analysis. Each axis we surveyed operates within its own silo. Retrieved precedents are never routed to a structuring module. Structured facts and extracted entities stop at the output layer; no argument graph constructor picks them up. Prediction models run without argument-level context, and their outputs are never fed into any counterfactual reasoning component. The closest the field has come to integration is coupling retrieval with generation~\cite{nigam-etal-2025-nyayarag} or with case-based reasoning~\cite{wiratunga-etal-2024-cbrrag}, but even those systems span only two of the five axes. We have come to think this fragmentation is no longer a reflection of the field's youth; it is a design choice the community has been making by default, and one that needs to be made deliberately instead.

At the task level, three gaps stand out. Argument graph construction for Indian legal judgments has no published dataset, no trained model, and no evaluation benchmark. It is, by any measure, the largest structural blind spot in the field. Counterfactual what-if analysis has not been attempted at all in the Indian context, even though it is arguably the capability most directly relevant to practising lawyers, the people these systems are nominally being built for. And evaluation remains systematically tilted toward Supreme Court data~\cite{malik-etal-2021-ildc, nigam-etal-2024-predex}. District and high courts, which together account for the overwhelming bulk of India's pending caseload, lack reliable benchmarks. LegalEval's annotation scope reflects the same imbalance~\cite{modi-etal-2023-semeval}.

The data-level picture is no better. No annotated resource captures argument-level structure, that is, the claim-evidence-statute-precedent linkages that argument graph construction would require. Multilingual annotations are absent entirely; Hindi and regional language passages embedded in English judgments remain invisible to every model we surveyed. The large-scale corpora that enable LLM fine-tuning, with NyayaAnumana chief among them~\cite{nigam-etal-2025-nyayaanumana}, trade annotation depth for volume. PredEx, by contrast, offers expert-level explanations but at a fraction of the scale~\cite{nigam-etal-2024-predex}. No single resource delivers both. We suspect that until one does, the field will continue producing systems that are either large-scale but shallow, or deep but narrow.

We should be clear about what this paper contributes and what it does not. It is a survey, not a system. We offer the first structured review of Indian legal NLP organised around five functional capability axes, map the existing literature to each axis, catalogue the datasets and evaluation methods currently in use, and identify what remains unsolved. Our primary claim is that individual capabilities have matured enough to make integration feasible, and that the absence of integration, rather than the immaturity of any one axis, is now the binding constraint on meaningful progress. We hope this survey serves as a useful reference for anyone aiming to build toward a unified Indian legal AI architecture. We make no claim that building one will be easy.

%% ================================================================
%%  SECTION 11 — Ethical Considerations
%% ================================================================
\section{Ethical Considerations}

Train a model on historical Indian court judgments and it will learn what those judgments contain, including their biases. Outcome disparities correlated with gender, caste, and geography are well documented in Indian legal history, and a prediction system optimised to reproduce historical outcome distributions will reproduce those inequities unless someone intervenes deliberately during training and evaluation~\cite{medvedeva-etal-2020-echr}. We believe bias auditing across demographic and geographic subgroups must be a non-negotiable prerequisite for deployment, not a post-launch afterthought.

Explainability, in a legal setting, is not a feature; it is a procedural obligation. Parties to litigation have the right to understand the basis of decisions that affect their legal interests, and that right does not evaporate the moment part of the decision process becomes automated~\cite{deakin-etal-2023-xai-law}. The models that predict best are often the ones that explain least, and a generated paragraph that reads fluently and cites the right statutes may still be legally incoherent. PredEx's methodology, grounding explanation evaluation in expert annotation rather than ROUGE scores~\cite{nigam-etal-2024-predex}, points in the right direction, but the harder lesson is one the community still needs to internalise: fluency is not soundness.

One distinction must be preserved. Everything we have surveyed here, retrieval systems, structuring modules, prediction models, and explanation generators, belongs in the category of decision-support tools, not autonomous adjudication. The most capable systems currently available still exhibit significant limitations in coverage and generalisation that rule out unsupervised deployment~\cite{nigam-etal-2025-nyayaanumana, mersha-etal-2024-xai}. The goal is augmentation, not automation, and we believe losing sight of that distinction is the fastest way to erode the procedural safeguards on which the rule of law depends.
